\documentclass[a4paper,12pt]{article}
\usepackage[T2A]{fontenc}
\usepackage[utf8]{inputenc}
\usepackage[russian]{babel}
\usepackage{amsmath,amssymb,mathtools}
\renewcommand{\sfdefault}{cmss}
\renewcommand{\familydefault}{\sfdefault}
\usepackage{icomma}
\usepackage[left=22mm,right=22mm,top=22mm,bottom=22mm]{geometry}
\usepackage{microtype}
\usepackage{tikz}
\usetikzlibrary{arrows.meta,positioning,shapes.geometric,calc}
\usepackage{booktabs,tabularx}
\usepackage{enumitem}
\usepackage{fancyhdr}
\usepackage{setspace}
\usepackage{xcolor}
\usepackage[hidelinks]{hyperref}
\definecolor{accent}{HTML}{126E71}
\definecolor{darkaccent}{HTML}{26494C}
\definecolor{textgray}{HTML}{455056}
\definecolor{lightline}{HTML}{CAD9D7}
\onehalfspacing
\setlength{\parindent}{0pt}
\setlength{\parskip}{4pt}
\emergencystretch=2em
\setlist[enumerate]{leftmargin=1.8em, itemsep=4pt, topsep=4pt}
\setlist[itemize]{leftmargin=1.6em, itemsep=2pt, topsep=2pt}
\pagestyle{fancy}
\fancyhf{}
\fancyhead[L]{\small\color{textgray}АННА БАННОВА\quad /\quad 16.10.26}
\fancyhead[R]{\small\color{textgray}МАТЕМАТИКА\quad /\quad ЕГЭ}
\fancyfoot[C]{\small\color{textgray}\thepage}
\renewcommand{\headrulewidth}{0.35pt}
\renewcommand{\footrulewidth}{0pt}
\setlength{\headheight}{15pt}
\newcommand{\sect}[1]{\vspace{6pt}{\large\bfseries\color{darkaccent}#1}\par\vspace{3pt}\noindent\color{lightline}\rule{\linewidth}{0.5pt}\color{black}\vspace{2pt}}
\newcommand{\keyphrase}[1]{{\color{accent}\bfseries #1}}
\newcommand{\unit}[1]{\,\text{#1}}
\begin{document}
\thispagestyle{empty}
\begin{center}
\vspace*{34mm}
{\small\bfseries\color{accent} МАТЕМАТИКА \quad / \quad ПРОФИЛЬНЫЙ ЕГЭ}\par
\vspace{15mm}
{\fontsize{31}{38}\selectfont\bfseries\color{darkaccent}Чек-лист}\par
\vspace{4mm}
{\fontsize{20}{27}\selectfont Степени и задачи\\прикладного характера}\par
\vspace{16mm}
{\color{lightline}\rule{94mm}{0.7pt}}\par
\vspace{13mm}
{\large Анна Баннова}\par
\vspace{3mm}
{\normalsize 16 октября 2026 года}\par
\vspace{25mm}
{\small\color{textgray}Повторение свойств степеней, решение степенных уравнений,\\модели силы натяжения и радиоактивного распада}\par
\vfill
{\small\bfseries\color{darkaccent}Лёвин Артём Александрович эксклюзивно для Анны Банновой}\par
\vspace{8mm}
\begin{tikzpicture}[x=1cm,y=1cm]
\draw[accent!55,line width=1.1pt] (-7,0) .. controls (-3.9,0.55) and (-1.4,-0.35) .. (0.5,0.15)
.. controls (2.7,0.60) and (4.5,-0.13) .. (7,0.22);
\end{tikzpicture}
\end{center}
\clearpage

\sect{01\quad Работа со степенями}
Для $a\ne0$ и любых допустимых показателей $r,s$:
\[
 a^r\cdot a^s=a^{r+s},\qquad
 \frac{a^r}{a^s}=a^{r-s},\qquad
 (a^r)^s=a^{rs}\quad\text{(если преобразование допустимо)}.
\]
Для дробных показателей всегда проверяйте, при каких $a$ определено выражение:
\[
 a^{p/q}=\bigl(\sqrt[q]{a}\bigr)^p,\quad q\in\mathbb N,\qquad
 a^{-r}=\frac1{a^r}\quad(a\ne0).
\]
Здесь дробь $p/q$ несократима. При чётном $q$ требуется $a\ge0$; при нечётном $q$ возможны и отрицательные $a$. Если показатель отрицательный, дополнительно $a\ne0$.

\keyphrase{Приём из повторения.} Разложите составные числа на множители, раскройте произведения степеней, а затем сложите показатели при одинаковых основаниях:
\[
\left(\frac35\right)^{1/8}\!\cdot5^{1/4}\cdot15^{7/8}
=3^{1/8+7/8}\cdot5^{-1/8+2/8+7/8}
=3\cdot5=15.
\]
\textbf{Отрицательная степень у дроби:}
\[
\left(\frac54\right)^{-1}=\frac45,
\qquad \frac{1}{\frac54}=\frac45,
\qquad \frac{5^{-1}}{4}=\frac1{20}.
\]
Скобки принципиальны: в последнем выражении показатель $-1$ относится только к числу $5$.

\sect{02\quad Как решать уравнения со степенями}
Главная цель --- выразить $x$ из степени, сохранив \keyphrase{все действительные корни}.

\textbf{Случай 1. Знаменатель дробного показателя чётный.}
\[
 x^{3/2}=64,\quad x\ge0
 \iff (\sqrt{x})^3=64
 \iff \sqrt{x}=4
 \iff \boxed{x=16}.
\]
\textbf{Случай 2. Знаменатель дробного показателя нечётный.}
\[
 x^{4/3}=81
 \iff (\sqrt[3]{x})^4=81
 \iff \sqrt[3]{x}=\pm3
 \iff \boxed{x=\pm27}.
\]
\textit{Почему два корня?} Чётная четвёртая степень даёт одно положительное значение как при $3$, так и при $-3$. Например, $(-8)^{4/3}=16$ и $8^{4/3}=16$.

\textbf{Случай 3. Отрицательный целый показатель.}
\[
 x^{-4}=81,\quad x\ne0
 \iff x^4=\frac1{81}
 \iff \boxed{x=\pm\frac13}.
\]
Если степень нечётная, ситуация иная: $x^{-3}=-\frac1{125}$ даёт только $x=-5$.

\keyphrase{Проверяйте себя.} Возведение в чётную степень стирает знак. Поэтому из $u^4=16$ следует $u=\pm2$, а из $u^3=8$ --- только $u=2$.

\sect{03\quad Как читать прикладную задачу}
\begin{enumerate}
\item Найдите в условии формулу и определите, за что отвечает каждый символ.
\item Выпишите известные величины \emph{с единицами измерения} и обозначьте искомую.
\item Сохраните знак условия: «не больше» означает $\le$, «не меньше» --- $\ge$.
\item Подставьте числа. При необходимости меняйте части равенства местами: $A=B\iff B=A$; при переносе одного слагаемого меняется знак.
\item Решите полученное степенное или показательное уравнение/неравенство, проверьте физические ограничения.
\item Если ищете наибольшее время или размер, исследуйте \emph{границу допустимых значений} и проверьте направление неравенства.
\end{enumerate}

\sect{04\quad Прикладная модель: кубический зонд}
Сила натяжения троса при подъёме зонда определяется формулой
\[
 F=\rho gL^3-mg,
\]
где $F$ --- сила натяжения (Н), $\rho$ --- плотность воды (кг/м$^3$), $g$ --- ускорение свободного падения (м/с$^2$), $L$ --- линейный размер кубического зонда (м), $m$ --- масса (кг).

\textbf{Пример занятия.} $F\le1000\unit{Н}$, $\rho=1000\unit{кг/м\textsuperscript{3}}$, $g=10\unit{м/с\textsuperscript{2}}$, $m=25\unit{кг}$. Найти наибольшее $L\ge0$.
\[
\begin{aligned}
1000\cdot10\cdot L^3-25\cdot10&\le1000,\\
10000L^3&\le1250,\\
L^3&\le\frac{125}{1000}=\left(\frac5{10}\right)^3,\\
\boxed{0\le L\le0{,}5\unit{м}}.&
\end{aligned}
\]
\textbf{Ответ:} максимальный размер $0{,}5\unit{м}$. При $L\ge0$ выражение $10000L^3-250$ возрастает; поэтому максимум достигается, когда сила достигает предельно допустимого значения.

\keyphrase{Как распознать точный куб:} $216=2^3\cdot3^3=(2\cdot3)^3=6^3$, отсюда
\[
\sqrt[3]{\frac{216}{1000}}=\frac6{10}=0{,}6.
\]
\sect{05\quad Прикладная модель: радиоактивный распад}
Масса изотопа через $t$ минут задаётся формулой
\[
 m(t)=m_0\cdot 2^{-t/T},
\]
где $m_0>0$ --- начальная масса, $T>0$ --- период полураспада, $t\ge0$ --- прошедшее время. Обозначение $m(t)$ означает массу \emph{в момент} $t$.

\textbf{Пример занятия.} $m_0=40\unit{мг}$, $T=10\unit{мин}$. Как долго масса будет \emph{не меньше} $5\unit{мг}$?
\[
\begin{aligned}
40\cdot2^{-t/10}&\ge5,\\
2^{-t/10}&\ge\frac18=2^{-3},\\
-\frac{t}{10}&\ge-3,\\
\boxed{0\le t\le30\unit{мин}}.&
\end{aligned}
\]
При умножении на отрицательное число знак неравенства меняется. Масса с течением времени убывает, поэтому через $30$ минут она достигнет порога, после чего станет меньше $5\unit{мг}$. \textbf{Ответ:} в течение первых $30$ минут.

\clearpage
\sect{Алгоритм\quad Прикладная задача со степенями}
\vspace{10mm}
\begin{center}
\begin{tikzpicture}[
 >={Stealth[length=3mm]}, line width=.9pt, font=\small,
 start/.style={ellipse, draw=accent, text width=5cm, align=center, minimum height=1.05cm, inner sep=5pt},
 action/.style={rectangle, rounded corners=2pt, draw=lightline!75!black, text width=6.2cm, align=center, minimum height=1.5cm, inner sep=7pt},
 decision/.style={diamond, aspect=2.1, draw=accent, align=center, text width=3.8cm, inner sep=3pt, minimum height=1.8cm},
 arr/.style={->,draw=darkaccent},
 label/.style={font=\footnotesize\bfseries,fill=white,inner sep=3pt,text=accent}
]
\node[start] (a) at (0,0) {Прочитать условие};
\node[action] (b) at (0,-2.1) {Записать формулу, величины, единицы и знак ограничения};
\node[decision] (d) at (0,-4.8) {Как неизвестная входит в формулу?};
\node[action] (left) at (-4,-7.7) {Через $L^3$: оставить куб слева и выразить $L^3$};
\node[action] (right) at (4,-7.7) {Через $2^{-t/T}$: привести обе части к основанию $2$};
\node[action] (left2) at (-4,-10.7) {Извлечь кубический корень; учесть $L\ge0$ и монотонность};
\node[action] (right2) at (4,-10.7) {Сравнить показатели; учесть знак при преобразовании неравенства};
\node[action] (c) at (0,-13.7) {Проверить подстановку, допустимость и нужную границу};
\node[start] (e) at (0,-15.9) {Записать ответ с единицей измерения};
\draw[arr] (a) -- (b);
\draw[arr] (b) -- (d);
\draw[arr] (d.west) -| node[label,near start,above] {$L^3$} (left.north);
\draw[arr] (d.east) -| node[label,near start,above] {$2^{-t/T}$} (right.north);
\draw[arr] (left) -- (left2);
\draw[arr] (right) -- (right2);
\draw[arr] (left2.south) |- ($(c.north)+(0,0.5)$) -- (c.north);
\draw[arr] (right2.south) |- ($(c.north)+(0,0.5)$);
\draw[arr] (c) -- (e);
\end{tikzpicture}
\end{center}

\clearpage
\sect{Самостоятельная тренировка\quad Путь А}
\textcolor{textgray}{10 заданий. Первые пять --- вычисления и уравнения, следующие пять --- применение степенных моделей. Во всех задачах выпишите ключевые преобразования.}
\begin{sloppypar}
\begin{enumerate}[label=\textbf{\arabic*.},itemsep=9pt,leftmargin=2em]
\item Вычислите: $\displaystyle\frac{2^{7/6}\cdot2^{1/3}}{2^{1/2}}$.
\item Упростите: $\displaystyle\left(\frac27\right)^{1/6}\cdot7^{1/3}\cdot14^{5/6}$.
\item Решите уравнение в действительных числах: $x^{3/2}=64$.
\item Решите уравнение в действительных числах: $x^{4/3}=81$.
\item Решите уравнение в действительных числах: $x^{-4}=16$.
\item Для кубического зонда $F=\rho gL^3-mg$. При $\rho=1000\unit{кг/м\textsuperscript{3}}$, $g=10\unit{м/с\textsuperscript{2}}$, $m=20\unit{кг}$ сила не должна превышать $1050\unit{Н}$. Найдите максимальный линейный размер $L$ (в метрах).
\item Для другого зонда действует та же формула. При $\rho=1000\unit{кг/м\textsuperscript{3}}$, $g=10\unit{м/с\textsuperscript{2}}$, $m=18\unit{кг}$ допустимая сила $F\le1980\unit{Н}$. Найдите максимальный $L$ (в метрах).
\item Масса изотопа изменяется по формуле $m(t)=m_0\cdot2^{-t/T}$. Вначале масса составляла $48\unit{мг}$, период полураспада $8\unit{мин}$. В течение скольких минут масса будет не меньше $6\unit{мг}$?
\item Масса изотопа описывается той же формулой. $m_0=96\unit{мг}$, $T=12\unit{мин}$. В течение скольких минут масса будет не меньше $3\unit{мг}$?
\item Начальная масса вещества $80\unit{мг}$, период полураспада $15\unit{мин}$. Через сколько минут \emph{впервые} масса станет не больше $10\unit{мг}$?
\end{enumerate}
\end{sloppypar}
\vfill
\textcolor{textgray}{Перед сдачей проверьте:}\par
\(\square\) учтены действительные корни;\quad
\(\square\) сохранён смысл «не больше / не меньше»;\par
\(\square\) числа представлены удобными степенями;\quad
\(\square\) ответы имеют единицы измерения.
\clearpage
\sect{Ответы для самопроверки}
\vspace{5mm}
\renewcommand{\arraystretch}{1.7}
\begin{tabularx}{\linewidth}{@{}>{\bfseries}p{12mm}X@{}}
\toprule
\textcolor{darkaccent}{№} & \textcolor{darkaccent}{Ответ} \\
\midrule
1 & $2$ \\
2 & $14$ \\
3 & $x=16$ \\
4 & $x=-27$ или $x=27$ \\
5 & $x=-\dfrac12$ или $x=\dfrac12$ \\
6 & $L_{\max}=0{,}5\unit{м}$ \\
7 & $L_{\max}=0{,}6\unit{м}$ \\
8 & $24\unit{мин}$ \\
9 & $60\unit{мин}$ \\
10 & $45\unit{мин}$ \\
\bottomrule
\end{tabularx}
\vspace{12mm}
\textbf{Особое внимание --- заданиям 4 и 5.} Если показатель при $x$ приводит к чётной степени неизвестной, необходимо проверить оба знака. В прикладных задачах подставьте найденную границу в исходную формулу и убедитесь, что неравенство выполняется именно на нужном промежутке.
\vfill
\begin{center}
\small\color{textgray}Лёвин Артём Александрович эксклюзивно для Анны Банновой
\end{center}
\end{document}
